% ============================================================================
% Rapport — Biais inductif de qualité d'image pour MedVAE (partie Théophile)
%
% NOTE POUR LA FUSION : contenu destiné à remplacer \section{Section 2:
% Théophile} (vide) dans final.tex, entre la section IV (JEPA) et la
% section VI (Segmentation). Le document maître charge déjà : graphicx,
% amsmath/amssymb, booktabs, cite, balance, float, dblfloatfix, subcaption.
%
% Copier le dossier figures/theophile/ à côté de final.tex (au même niveau
% que figures/elias/, figures/yanic/, figures/theo/) avant de coller le
% contenu ci-dessous (entre BEGIN/END SECTION) dans final.tex.
%
% Ce fichier compile aussi seul (préambule ci-dessous, à supprimer lors de
% la fusion) pour prévisualisation.
% ============================================================================
\documentclass[conference]{IEEEtran}
\usepackage[utf8]{inputenc}
\usepackage[T1]{fontenc}
\usepackage{graphicx}
\usepackage{amsmath,amssymb}
\usepackage{booktabs}
\usepackage{cite}
\usepackage{balance}
\usepackage{float}
\usepackage{dblfloatfix}
\usepackage{subcaption}
\usepackage{placeins}

\begin{document}
\pagestyle{plain}

% ============================================================
% BEGIN SECTION — à coller à la place de \section{Section 2: Théophile}
% ============================================================
\section{Biais inductif de qualité d'image pour MedVAE}
\label{sec:theophile}

\subsection{Motivation}

La section précédente a montré que MedVAE se comporte comme un filtre
passe-bas intelligent et un débruiteur non linéaire, mais qu'il ne dispose
d'aucun moyen explicite de savoir \emph{à quel point} l'image qu'on lui
présente est dégradée : il applique la même transformation à une
angiographie propre ou fortement bruitée. Or cette information est
disponible \emph{a priori} et peut être estimée par des métriques de
qualité sans référence (NR-IQA). Nous formulons l'hypothèse que
conditionner explicitement l'autoencodeur par un score de qualité
$c \in [0,1]$ améliore sa reconstruction, en lui permettant d'adapter son
comportement au niveau de dégradation réellement présent en entrée. Trois
questions en découlent : comment calculer $c$ de façon objective et
reproductible en l'absence d'annotation humaine sur ARCADE ; comment
l'injecter dans MedVAE sans réentraîner le modèle depuis zéro ; et si le
gain observé, le cas échéant, est attribuable au contenu informatif de $c$
ou à un simple effet de capacité supplémentaire. Cette dernière question est
loin d'être anodine : tout mécanisme de conditionnement ajoute des
paramètres et change la dynamique d'entraînement, de sorte qu'un gain
apparent pourrait simplement refléter une capacité accrue plutôt qu'une
réelle exploitation du signal de qualité, d'où la nécessité d'une ablation
dédiée (section~\ref{sec:ablation}). Pour la première question, trois
méthodes indépendantes de calcul de $c$ (approches A, B, C)
sont comparées en parallèle, et l'ensemble du pipeline en aval
(conditionnement, entraînement, évaluation, ablation) est exécuté
séparément pour chacune,
ce qui permet de distinguer un effet propre à une méthode de calcul de $c$
d'un effet plus général du conditionnement FiLM lui-même.

\subsection{Construction du score de qualité}
\label{sec:score}

Sur le sous-ensemble segmentation d'ARCADE (1000 images train / 200
validation, $512\times512$, niveaux de gris), chaque image est découpée en
une grille de $4\times4$ patchs, une granularité locale motivée par le
caractère spatialement non uniforme des artefacts en imagerie cardiaque
(mouvement, bruit) \cite{sekiguchi2025patchdsa}, normalisée par \emph{Mean
Subtracted Contrast Normalization} ($\hat{I}=(I-\mu)/(\sigma+C)$, fenêtre
$7\times7$), puis caractérisée par 8 métriques par patch : variance du
Laplacien, contraste RMS, entropie de Shannon, et le gradient de Tenengrad
décliné en 4 directions, une analyse directionnelle adaptée à la structure
tubulaire des vaisseaux \cite{wang2024noreference} (ce qui permet de
détecter le flou de mouvement cardiaque, anisotrope). Ces métriques sont
agrégées en 36 features par image. Un
ResNet-18 pré-entraîné (poids figés) extrait en parallèle un vecteur de 512
features, utilisé uniquement par l'approche apprise : les deux branches de
ce pipeline de features (Figure~\ref{fig:featpipe}) ne convergent jamais
directement, elles alimentent simplement des méthodes de calcul du score
différentes (section~\ref{sec:score}).

\begin{figure}[t]
  \centering
  \includegraphics[width=0.85\columnwidth]{pipeline_A_features.pdf}
  \caption{Construction des features de qualité : métriques IQA par patch
  et features ResNet-18.}
  \label{fig:featpipe}
\end{figure}

Ces features sont ensuite fusionnées en un score unique $c$, $1$
correspondant à une image cliniquement exploitable, selon trois méthodes
indépendantes (Figure~\ref{fig:scorepipe}), toutes fittées exclusivement
sur le sous-ensemble TRAIN. L'approche B applique une PCA (après
standardisation) sur 6 features sélectionnées (les 36 précédentes plus deux
features d'homogénéité spatiale et d'équilibre directionnel) ; le premier
axe explique 56{,}1\,\% de la variance, et son signe est corrigé par
corrélation avec le Tenengrad ($r=0{,}952$). L'approche A pondère
linéairement les mêmes 6 features par les \emph{loadings} absolus
normalisés de ce même PC1. L'approche C, faute d'annotation humaine sur
ARCADE, apprend un score sur un \emph{proxy synthétique}, une stratégie de
supervision indirecte déjà employée pour calibrer des scores NR-IQA sur
des images biomédicales sans vérité terrain perceptuelle
\cite{asri2026deeplearning} : chaque image TRAIN est dégradée
artificiellement (bruit, flou, JPEG) avec une cible
$y=1-\text{sévérité connue}$, et un MLP léger ($512{\to}128{\to}32{\to}1$)
est entraîné par MSE sur les features ResNet-18 de ce jeu synthétique (MSE
de validation $1{,}58\times10^{-3}$), puis appliqué aux images réelles.

\begin{figure*}[t]
  \centering
  \includegraphics[width=0.92\textwidth]{pipeline_B_score.pdf}
  \caption{Les trois approches de calcul du score de qualité $c$.}
  \label{fig:scorepipe}
\end{figure*}

A et B, construites sur les mêmes métriques, sont fortement corrélées
($r=0{,}884$) ; C en est en revanche quasiment indépendante
($r(A,C)=0{,}003$, $r(B,C)=0{,}015$, Table~\ref{tab:corr}) et sa
distribution est beaucoup plus resserrée vers 1 (moyenne 0{,}890) que A
(0{,}448) ou B (0{,}483) : le MLP appris sur des features profondes capture
une notion de qualité différente d'une combinaison de métriques classiques.
Le score actif est enfin discrétisé en trois classes
\emph{bad}/\emph{medium}/\emph{good} par tertiles (TRAIN). La
Figure~\ref{fig:gallery} illustre cette progression sur des images réelles
triées par score croissant : les angiographies les moins nettes et les plus
bruitées (gauche) cèdent progressivement la place à des acquisitions
contrastées où l'arbre vasculaire est bien défini (droite), ce qui confirme
qualitativement que le score capture une notion de qualité d'image
cohérente avec une lecture clinique.

\begin{figure}[t]
  \centering
  \includegraphics[width=\columnwidth]{result_gallery_score.png}
  \caption{Galerie d'images ARCADE triées par score de qualité croissant
  (approche A, pondéré) : de la moins nette (gauche) à la plus contrastée
  (droite).}
  \label{fig:gallery}
\end{figure}

\begin{table}[htbp]
  \centering
  \begin{tabular}{@{}lccc@{}}
    \toprule
     & A (pondéré) & B (PCA) & C (appris) \\
    \midrule
    A (pondéré) & 1 & 0{,}884 & 0{,}003 \\
    B (PCA)     &   & 1       & 0{,}015 \\
    C (appris)  &   &         & 1 \\
    \bottomrule
  \end{tabular}
  \caption{Corrélation de Pearson entre les trois scores (1200 images).}
  \label{tab:corr}
\end{table}

\FloatBarrier
\subsection{Conditionnement FiLM et entraînement}
\label{sec:film}

Pour conditionner MedVAE \cite{varma2025medvae} (\texttt{AutoencoderKL}
pré-entraîné \cite{kingma2014autoencoding}, variante
\texttt{medvae\_4\_1\_2d}, 57{,}06M paramètres) par $c$ sans modifier sa
topologie d'entrée, nous utilisons la \emph{Feature-wise Linear Modulation}
(FiLM \cite{perez2018film}) plutôt qu'une concaténation de canal : un MLP
partagé encode $c$ en un embedding de 128 dimensions, et 19
têtes linéaires (une par \texttt{ResnetBlock} de l'encodeur, du goulot et
du décodeur) produisent des paramètres affines appliqués par un hook sur
chaque bloc, $h \leftarrow (1+\gamma_i(c)) \cdot h + \beta_i(c)$, avec
initialisation à zéro (le modèle conditionné est donc identique au modèle
original avant tout entraînement). Ce choix d'initialisation est important :
il garantit que toute différence de performance observée après
entraînement provient de ce que le réseau a appris à faire de $c$, et non
d'une perturbation aléatoire de poids préentraînés. FiLM est par ailleurs
préféré à une concaténation de canal d'entrée car il laisse la topologie de
l'encodeur intacte et concentre l'effet du conditionnement sur des
paramètres clairement identifiables, ce qui facilite l'ablation de la
section~\ref{sec:ablation}. Ce mécanisme ajoute 1{,}70M paramètres
(2{,}98\,\%). L'entraînement se déroule en deux phases sur les 1000 images
TRAIN redimensionnées à $64\times64$ : un warm-up où seul FiLM est
entraîné (2500 steps, taux d'apprentissage faible) pour éviter de
déstabiliser le décodeur préentraîné avec des paramètres affines encore
mal calibrés, puis un fine-tuning global où tous les poids sont dégelés
(5000 steps, \emph{cosine annealing}) pour permettre au reste du réseau de
s'adapter à la modulation FiLM. Un VAE baseline strictement identique mais
sans FiLM est entraîné en parallèle sur le même budget de
7500 steps, afin d'isoler l'effet propre du conditionnement de celui d'un
fine-tuning supplémentaire (Figure~\ref{fig:filmpipe}).

\begin{figure*}[t]
  \centering
  \includegraphics[width=0.92\textwidth]{pipeline_C_film.pdf}
  \caption{Conditionnement FiLM (19 blocs) et entraînement en deux phases,
  avec ablation baseline.}
  \label{fig:filmpipe}
\end{figure*}

\begin{table}[htbp]
  \centering
  \begin{tabular}{@{}lccc@{}}
    \toprule
    Approche & cVAE($c$) & Baseline & Gain relatif \\
    \midrule
    A (pondéré) & 0{,}00735 & 0{,}01345 & $-45{,}3\%$ \\
    B (PCA)     & 0{,}00749 & 0{,}01348 & $-44{,}5\%$ \\
    C (appris)  & 0{,}00925 & 0{,}01345 & $-31{,}2\%$ \\
    \bottomrule
  \end{tabular}
  \caption{Perte ELBO de validation finale ($64\times64$, interne à
  l'entraînement), cVAE conditionné vs baseline, par approche.}
  \label{tab:trainloss}
\end{table}

Sur cette perte interne, le cVAE conditionné est nettement meilleur que le
baseline pour les trois approches (Table~\ref{tab:trainloss}), un résultat
qui semble valider l'hypothèse de départ, mais qui ne mesure que
l'optimisation sur le format réduit du fine-tuning, pas la fidélité de
reconstruction sur des images de test indépendantes en pleine résolution.

\FloatBarrier
\subsection{Évaluation de la fidélité de reconstruction}
\label{sec:eval}

Sur un jeu de test de 60 images (20 par classe, non vues à l'entraînement),
nous comparons la reconstruction du cVAE et du baseline selon quatre
métriques (MSE, SSIM, PSNR, HaarPSI, cette dernière étant recommandée pour
les images médicales \cite{breger2024adequacy}), à la fois sur l'image
entière et sur un PSNR masqué restreint aux pixels de vaisseaux (mPSNR,
annotations COCO rastérisées). Cette restriction permet de vérifier qu'un
éventuel résultat n'est pas un simple artefact du fond, non pertinent
cliniquement et potentiellement trompeur puisqu'un PSNR global peut
\emph{augmenter} quand l'image se dégrade
% TODO fusion : remplacer par \ref{...} vers la section robustesse de MedVAE
(cf. l'analyse de robustesse de MedVAE présentée précédemment)
(Figure~\ref{fig:evalpipe}).

\begin{figure}[t]
  \centering
  \includegraphics[width=0.82\columnwidth]{pipeline_D_evaluation.pdf}
  \caption{Protocole d'évaluation : baseline vs cVAE, PSNR classique et
  masqué (vaisseaux).}
  \label{fig:evalpipe}
\end{figure}

\begin{table}[htbp]
  \centering
  \begin{tabular}{@{}lcccc@{}}
    \toprule
    Approche & MSE & SSIM & PSNR & HaarPSI \\
    \midrule
    A (pondéré) & $+21{,}1\%$ & $-1{,}6\%$ & $-2{,}4\%$ & $-1{,}6\%$ \\
    B (PCA)     & $+26{,}1\%$ & $-1{,}7\%$ & $-2{,}7\%$ & $-1{,}9\%$ \\
    C (appris)  & $+48{,}8\%$ & $-1{,}8\%$ & $-4{,}4\%$ & $-2{,}1\%$ \\
    \bottomrule
  \end{tabular}
  \caption{Variation relative du cVAE par rapport au baseline (60 images
  réelles ; MSE : positif = pire ; les autres : négatif = pire). Toutes les
  différences sont significatives (Wilcoxon $p<0{,}0001$).}
  \label{tab:eval}
\end{table}

Le résultat est contraire à l'hypothèse de départ : sur les trois approches
et les quatre métriques, sans exception, le baseline non conditionné
reconstruit mieux que le cVAE (Table~\ref{tab:eval}), et cela
ne change pas en restreignant l'erreur aux vaisseaux : le mPSNR confirme le
même classement pour les trois approches, avec un écart systématique de
$2$ à $5$~dB par rapport au PSNR classique (les vaisseaux, fins et peu
contrastés, restent intrinsèquement plus difficiles à reconstruire que le
fond,
% TODO fusion : remplacer par \ref{...} vers la section robustesse de MedVAE
ce qui est cohérent avec l'analyse de robustesse de MedVAE présentée
précédemment). Le détail par classe (Table~\ref{tab:mpsnr-class}) confirme
que ce classement baseline $>$ cVAE est uniforme sur les 9 combinaisons
approche$\times$classe (bad/medium/good), sans aucune exception, et se
confirme sur des dégradations synthétiques contrôlées (bruit, flou de
mouvement) : le baseline égale ou dépasse le cVAE sur la quasi-totalité des
niveaux testés. La perte de validation interne (Table~\ref{tab:trainloss}),
plus favorable au cVAE, s'avère donc
un indicateur trompeur de la performance réelle.

\begin{table}[htbp]
  \centering
  \footnotesize
  \begin{tabular}{@{}lcccccc@{}}
    \toprule
    & \multicolumn{2}{c}{A (pondéré)} & \multicolumn{2}{c}{B (PCA)} & \multicolumn{2}{c}{C (appris)} \\
    Classe & base & cVAE & base & cVAE & base & cVAE \\
    \midrule
    bad    & 31{,}42 & 29{,}96 & 32{,}07 & 30{,}55 & 30{,}57 & 29{,}42 \\
    medium & 30{,}62 & 29{,}72 & 31{,}33 & 29{,}57 & 30{,}87 & 29{,}34 \\
    good   & 30{,}97 & 29{,}89 & 29{,}72 & 28{,}48 & 31{,}37 & 29{,}75 \\
    \bottomrule
  \end{tabular}
  \caption{mPSNR (dB) baseline vs cVAE, par approche et par classe de
  qualité (60 images de test). Le baseline gagne sur les 9 combinaisons,
  sans aucune exception.}
  \label{tab:mpsnr-class}
\end{table}

Cette régularité tient aussi à l'intérieur de chaque classe : la
dispersion de l'écart baseline/cVAE reste comparable entre les deux
modèles pour les trois métriques, ce n'est donc pas un effet dominé par
quelques images aberrantes, mais un décalage systématique de toute la
distribution.

\FloatBarrier
\subsection{Ablation par signal constant}
\label{sec:ablation}

Le résultat précédent invite à vérifier si le réseau exploite malgré tout
le contenu informatif de $c$, même sans en tirer un gain net. Un second
cVAE est entraîné avec exactement la même architecture, les mêmes données
et le même budget, mais en présentant $c=0{,}5$ constant à toutes les
images ; à l'inférence, les deux modèles reçoivent le vrai score de chaque
image de test, et seule la valeur \emph{vue à l'entraînement} diffère entre
les deux (Figure~\ref{fig:ablationpipe}).

\begin{figure}[t]
  \centering
  \includegraphics[width=0.78\columnwidth]{pipeline_E_ablation.pdf}
  \caption{Ablation par signal $c$ constant : même architecture, mêmes
  données, seule la supervision de $c$ change.}
  \label{fig:ablationpipe}
\end{figure}

\begin{table}[htbp]
  \centering
  \begin{tabular}{@{}lccc@{}}
    \toprule
     & $c$ réel gagne & $c{=}0{,}5$ gagne & Non sig. \\
    \midrule
    A (pondéré) & 0/4 & 2/4 & 2/4 \\
    B (PCA)     & 2/4 & 0/4 & 2/4 \\
    C (appris)  & 4/4 & 0/4 & 0/4 \\
    \bottomrule
  \end{tabular}
  \caption{Nombre de métriques (sur 4) où cVAE($c$ réel) bat
  significativement cVAE($c{=}0{,}5$) (Wilcoxon $p<0{,}05$, 60 images).}
  \label{tab:ablation}
\end{table}

Ce résultat, plus nuancé, dépend fortement de l'approche
(Table~\ref{tab:ablation}) : seule C démontre sans ambiguïté que le signal
$c$ est exploité, sur les quatre métriques ($p<10^{-4}$), constat qui se
confirme également en mPSNR sur les trois classes sans exception. Pour B,
l'effet est partiel (MSE et PSNR significatifs, pas SSIM/HaarPSI), et le
mPSNR par classe montre des inversions de verdict selon la métrique. Pour
A, l'effet est même inversé sur deux métriques perceptuelles :
$c=0{,}5$ constant y fait mieux que le vrai score pondéré.

La Figure~\ref{fig:delta} détaille, pour l'approche C, l'écart par image
entre cVAE($c$ réel) et cVAE($c{=}0{,}5$) en fonction du score de qualité
de l'image testée : l'avantage du conditionnement réel croît avec $c$, et
s'inverse légèrement pour les images les plus dégradées (\emph{bad}), où
$c=0{,}5$ constant, proche de la médiane, agit presque comme un signal par
défaut raisonnable. Le réseau a donc bien appris une réponse
\emph{dépendante} de $c$, et pas seulement une réponse à la présence d'un
embedding non nul.

\begin{figure}[t]
  \centering
  \includegraphics[width=\columnwidth]{result_delta_ablation.png}
  \caption{Écart par image (MSE, SSIM, HaarPSI) entre cVAE($c$ réel) et
  cVAE($c{=}0{,}5$) en fonction du score de qualité, approche C.}
  \label{fig:delta}
\end{figure}

\FloatBarrier
\subsection{Discussion}
\label{sec:discussion-theophile}

L'ensemble de ces expériences raconte une histoire plus riche qu'une simple
validation de l'hypothèse de départ. Le conditionnement FiLM, mesuré
correctement sur un jeu de test indépendant et en espace pixel, dégrade
légèrement la fidélité de reconstruction par rapport à un VAE non
conditionné entraîné avec le même budget. Ce résultat est robuste,
significatif sur trois approches et quatre métriques, et résiste à la
restriction aux pixels vasculaires. L'ablation nuance ce constat négatif :
pour l'approche C, le réseau exploite sans ambiguïté le contenu informatif
de $c$, donc la dégradation observée n'est pas due à un conditionnement
inerte, mais plutôt à un compromis architectural réel, où la capacité FiLM
supplémentaire et la modulation par couche semblent perturber
l'optimisation du décodeur plus qu'elles ne l'enrichissent, dans la limite
du budget utilisé. Pour A et B, l'effet de $c$ est plus faible voire
inversé, suggérant que leurs scores, des combinaisons linéaires de
métriques de netteté et de contraste, apportent une information déjà
largement accessible au réseau via ses propres convolutions.

La quasi-indépendance entre le score C et les scores A/B
(Table~\ref{tab:corr}, $r\approx0$) invite à la prudence sur
l'interprétation du score appris : entraîné sur un proxy synthétique de
sévérité plutôt qu'une vérité terrain humaine, il pourrait encoder un
signal partiellement différent de la qualité perceptuelle réelle. Que ce
signal soit exploité par le réseau sans pour autant améliorer la
reconstruction est cohérent avec cette lecture : le réseau apprend
\emph{quelque chose} d'utile à partir de $c_C$, mais pas nécessairement une
meilleure stratégie de reconstruction.

Plusieurs limites de ce protocole méritent d'être soulignées. D'abord, le
budget d'entraînement (7500 steps, résolution $64\times64$) reste modeste
au regard de la taille de MedVAE (57M paramètres) : un compromis
architectural FiLM qui pénalise le décodeur à ce stade pourrait s'atténuer,
voire s'inverser, avec un budget plus long ou un \emph{warm-up} plus
progressif. Ensuite, l'absence d'annotation humaine de qualité sur ARCADE
nous a contraints à valider indirectement les trois scores (cohérence
mutuelle, corrélation avec des métriques de référence) plutôt que contre
une vérité terrain perceptuelle ; une étude avec des cliniciens permettrait
de trancher si le score appris (C), le seul réellement exploité par le
réseau, capture une notion de qualité cliniquement pertinente ou un
artefact du proxy synthétique utilisé pour le superviser. Enfin, le critère
utilisé ici (fidélité pixel/structure) ne mesure pas directement l'utilité
clinique en aval (p.\,ex. lisibilité du réseau vasculaire pour un
diagnostic) ; un travail futur pourrait coupler ce protocole à une tâche de
segmentation ou de détection pour vérifier si le conditionnement FiLM, même
légèrement défavorable en reconstruction brute, n'apporte pas un bénéfice
sur une métrique orientée tâche.

% --------------------------------------------------------------------------
% NOTE POUR LA FUSION : ces références sont propres à cette partie. À
% fusionner manuellement avec la bibliographie commune de final.tex (qui
% liste déjà MedVAE et Kingma & Welling sous d'autres clés) : il faudra
% dédupliquer varma2025medvae/kingma2014autoencoding avec les entrées
% existantes et renuméroter les \cite en conséquence.
% --------------------------------------------------------------------------
\begin{thebibliography}{99}

\bibitem{perez2018film}
E.~Perez, F.~Strub, H.~de Vries, V.~Dumoulin, and A.~Courville,
``FiLM: Visual reasoning with a general conditioning layer,'' in
\emph{Proc. AAAI Conf. on Artificial Intelligence}, 2018.

\bibitem{breger2024adequacy}
A.~Breger \emph{et al.}, ``A study on the adequacy of common IQA measures
for medical images,'' \emph{arXiv:2405.19224}, 2024.

\bibitem{wang2024noreference}
Z.~Wang, S.~Chen, J.~Dai, S.~Qin, Y.~Cao, R.~Zhao, G.~Wu, Y.~Tang, and
J.~Chen, ``A no-reference medical image quality assessment method based on
automated distortion recognition technology: application to preprocessing
in MRI-guided radiotherapy,'' \emph{arXiv:2412.06599}, 2024.

\bibitem{asri2026deeplearning}
M.~A.~A.~Asri, H.~Rajagopal, N.~Mokhtar, W.~A.~W.~M.~Mahiyiddin,
Y.~A.~L.~Lim, M.~Iwahashi, R.~Harakawa, F.~Ibrahim, and T.~Ito,
``Deep learning-based no-reference image quality assessment framework for
\emph{Cryptosporidium} spp. and \emph{Giardia} spp.,'' \emph{PLOS ONE},
2026.

\bibitem{sekiguchi2025patchdsa}
Y.~Sekiguchi, T.~Okamoto, T.~Matsuzawa, K.~Fujimoto, K.~Fujiwara,
T.~Kondo, J.~Koizumi, and H.~Haneishi, ``PatchDSA: improving digital
subtraction angiography with patch-based phase-matching in natural
breathing scenarios,'' \emph{Radiological Physics and Technology}, vol.~18,
no.~3, pp.~698--706, 2025.

\bibitem{varma2025medvae}
M.~Varma, A.~Kumar, R.~van der Sluijs, S.~Ostmeier, L.~Blankemeier,
P.~Chambon, C.~Bluethgen, J.~Prince, C.~Langlotz, and A.~Chaudhari,
``MedVAE: efficient automated interpretation of medical images with
large-scale generalizable autoencoders,'' \emph{arXiv:2502.14753}, 2025.

\bibitem{kingma2014autoencoding}
D.~P.~Kingma and M.~Welling, ``Auto-encoding variational Bayes,'' in
\emph{Proc. Int. Conf. on Learning Representations (ICLR)}, 2014.

\end{thebibliography}

% ============================================================
% END SECTION
% ============================================================

\end{document}